\section{Adjacency Manipulation}

\subsection{The Attack}

The adjacency graph determines which tracts can be in the same district.
Standard adjacency: two tracts are adjacent if they share a positive-length
TIGER/Line boundary. An adversary could use a non-standard criterion, such as:
\begin{itemize}
  \item Minimum shared boundary length (e.g., require $\geq 100$m of shared boundary)
  \item Rook adjacency (shared boundary must be a line, not a point)
  \item Water adjacency exclusion (no adjacency across water bodies)
\end{itemize}

\subsection{Synthetic Poisoned Adjacency Test}

We test the maximum impact of adjacency manipulation by constructing a synthetic
``poisoned'' NC adjacency graph: we identify the 50 tract-pair edges that most
favour keeping Democratic and Republican tracts separated (edges where one tract
has $v_i > 0.6$ and the adjacent tract has $v_j < 0.4$), and remove them from
the adjacency graph.
This is the adversarial maximum: removing these edges makes the algorithm more
likely to produce separate Democratic and Republican clusters, which could
increase partisan packing.

Running \textsc{bisect} on the poisoned graph for NC:
\begin{table}[h]
\centering\small
\caption{NC adjacency manipulation results: standard vs.\ poisoned adjacency.
Partisan shift from removing 50 cross-partisan edges.}
\label{tab:adjacency-gaming}
\begin{tabular}{lccc}
\toprule
Adjacency & Democratic seats & PP compactness & $\Delta D$ \\
\midrule
Standard (TIGER/Line) & 7.0 & 0.171 & baseline \\
Poisoned ($-50$ edges) & 6.85 & 0.168 & $-0.15$ \\
\bottomrule
\end{tabular}
\end{table}

The poisoned adjacency produces a partisan shift of $-0.15$ Democratic seats ---
roughly comparable to the resolution gaming maximum.
The small magnitude confirms B.7's finding that seed and topology variation
have limited partisan effects.

\subsection{Audit Protection}

The SHA-256 hash of the adjacency graph (\texttt{adjacency\_hash} in the manifest)
detects any modification: the poisoned graph has a different hash from the standard
TIGER/Line adjacency.
The verifier can independently reconstruct the standard adjacency from the public
TIGER/Line files and compare SHA-256 hashes.

\begin{verbatim}
bisect label-verify --check-adjacency official_2020
\end{verbatim}

This command reconstructs the expected adjacency hash from TIGER/Line and verifies
it matches the manifest. Any adjacency manipulation is immediately detected.
